\documentclass{article}



% Language setting

% Replace `english' with e.g. `spanish' to change the document language

\usepackage[english]{babel}



% Set page size and margins

% Replace `letterpaper' with `a4paper' for UK/EU standard size

\usepackage[letterpaper,top=2cm,bottom=2cm,left=3cm,right=3cm,marginparwidth=1.75cm]{geometry}



% Useful packages

\usepackage{amsmath}

\usepackage{amssymb}

\usepackage{graphicx}

\usepackage[colorlinks=true, allcolors=black]{hyperref}

\usepackage{titlesec}
\usepackage{xparse}



\NewDocumentCommand{\qcq}{m o o m}{%
	\ensuremath{\forall #1 \IfValueT{#2}{, #2} \IfValueT{#3}{, #3}  \in  \mathbb{#4}}%
}
\NewDocumentCommand{\ext}{m o o m}{%
	\ensuremath{\exists #1 \IfValueT{#2}{, #2} \IfValueT{#3}{, #3}  \in  \mathbb{#4}}%
}



\title{\textbf{Exercices sur la continuité et la dérivabilité}}

\author{\textbf{Yassin Akkab}}



\newcommand{\R}{\mathbb{R}}
\newcommand{\N}{\mathbb{N}}
\newcommand{\Z}{\mathbb{Z}}
\newcommand{\C}{\mathbb{C}}
\newcommand{\Q}{\mathbb{Q}}
\newcommand{\D}{\mathbb{D}}





\graphicspath{ {./Screenshots/} }




\begin{document}
	
	\maketitle
	

	
	\subsection*{Exercice 1}
	Soit la fonction \( f(x) = 
	\begin{cases}
		2x + 1 & \text{si } x \leq 1, \\
		x^2 & \text{si } x > 1.
	\end{cases}
	\)
	\begin{itemize}
		\item \( f \) est-elle continue en \( x = 1 \) ?
		\item Si non, modifiez la fonction pour qu'elle devienne continue en \( x = 1 \).
	\end{itemize}
	
	\subsection*{Exercice 2}
	Soit la fonction \( f(x) = \frac{1}{x} \). Étudiez la continuité de cette fonction sur son domaine de définition.
	
	\subsection*{Exercice 3}
	Soit la fonction \( g(x) = |x| \). 
	\begin{itemize}
		\item La fonction \( g \) est-elle dérivable en \( x = 0 \) ?
	
	\end{itemize}
	
	\subsection*{Exercice 4}
	Calculez les dérivées des fonctions suivantes :
	\[
	f(x) = x^3 + 2x^2 - 5x + 7
	\]
	\[
	g(x) = \frac{1}{x^2}
	\]
	\[
	h(x) = \sqrt{x}
	\]
	
	\subsection*{Exercice 5}
	Soit la fonction \( f(x) = e^x \). 
	\begin{itemize}
		\item Montrez que \( f \) est dérivable sur \( \mathbb{R} \).
		\item Calculez \( f'(x) \).
	\end{itemize}
	
	\subsection*{Exercice 6}
	Étudiez la continuité et la dérivabilité de la fonction suivante :
	\[
	f(x) = \frac{x^2 - 1}{x - 1}
	\]
	pour \( x \neq 1 \).
	
\subsection*{Exercice 7 : Étude de fonctions usuelles}
Étudiez les fonctions suivantes :
\begin{itemize}
	\item \( f_1(x) = x \)
	\item \( f_2(x) = x^2 \)
	\item \( f_3(x) = \frac{1}{x} \)
	\item \( f_4(x) = \sqrt{x} \)
	\item \( f_5(x) = \sin(x) \)
	\item \( f_7(x) = \tan(x) \)
\end{itemize}



	
	
	

	
	
	
	
	
	
	
	
\end{document}